\begin{theorem}
        Let
        $\RangeSpace_1 = \pth{\GroundSet,\RangeSet^1}, \ldots,
        \RangeSpace_k = \pth{\GroundSet,\RangeSet^k}$
        be range spaces with \VC dimension $\Dim_1, \ldots, \Dim_k$,
        respectively.  Next, let $f\pth{ \range_1, \ldots, \range_k}$
        be a function that maps any $k$-tuple of sets
        $\range_1 \in \RangeSet^1, \ldots, \range_k \in \RangeSet^k$
        into a subset of $\GroundSet$. Here, the function $f$ is
        restricted to be defined by a sequence of set operations like
        complement, intersection and union.  Consider the range set
        \begin{align*}
          \RangeSetA%
          =%
          \Set{ f\pth{ \range_1, \ldots, \range_k}}{ \range_1
          \in \RangeSet_1, \ldots, \range_k \in \RangeSet_k \bigr.}
        \end{align*}
        and the associated range space $\RangeSpaceA =
        \pth{\GroundSet, \RangeSetA}$. Then, the \VC dimension of
        $\;\RangeSpaceA$ is bounded by $O\pth{ k \Dim \lg k}$, where
        $\Dim = \max_i \Dim_i$.
        
        \thmlab{VC:combination}
    \end{theorem}
